\begin{frame}{기원: ADALINE의 ``자격 있는 시냅스''}
  \begin{itemize}
    \item 이름이 생소해 보이지만 기원은 \textbf{1980년대 초} ---
      Barto, Sutton, Anderson (1983$\sim$89) 논문.
    \item ADALINE(최소제곱학습기)에 ``그 시점에 학습에
      \textbf{자격(eligibility)}이 있는 시냅스가 어디인가''를 기록하는
      장치를 처음 도입.
    \item 즉, 강화학습 TD($\lambda$)가 아니라
      \textbf{``어떤 시냅스가 지금 보상을 받을 자격이 있는가''를
      묻던 신경망의 부가장치}였던 셈.
    \item ``어떤 \textbf{상태}가 지금 TD 오차에 기여할 자격이 있는가''
      로 옮겨온 것이 지금의 적격흔적.
  \end{itemize}
  \vskip0.8em
  \kq{``자격을 남기고, 나중에 그 자격만큼 크레딧을 배분한다''.}
\end{frame}
